% Source: https://artofproblemsolving.com/community/c3h1647789_difficult_geometrycp_problem
% Stable identity: aops-topic-1647789-q01
Call the jumps $U:(x,y)\mapsto(x-4,y+6)$ and $V:(x,y)\mapsto(x+1,y-1)$. $U$ is allowed only when $x\ge4$, and $V$ only when $y\ge1$.

\textbf{Escaping.} Four $V$ jumps and one $U$ jump move the flea by $(0,2)$. If the flea is at a point with $x\ge4$, it can jump $U$ to $(x-4,y+6)$; since now $y+6\ge4$, it can jump $V$ four times to $(x,y+2)$, and then $x\ge4$ again. Repeating, it climbs by $2$ each round, so it gets as far from the origin as it likes, in particular farther than $2018$. So every point from which the flea can reach a point with $x\ge4$ is ``bad'' (not in the set).

\textbf{Trapped points.} Suppose $x<4$. The only possible jump is $V$ until the flea first reaches $x\ge4$. Let $k$ be the least positive integer with $x+k\ge4$ (so $k=4,3,2,1$ for $x$ in $[0,1)$, $[1,2)$, $[2,3)$, $[3,4)$). The flea can make $k$ jumps $V$ exactly when $y\ge k$. Hence:
\begin{itemize}
\item if $y\ge k$, the flea reaches $(x+k,y-k)$ with $x+k\ge4$ and escapes as above;
\item if $y<k$, it can make at most $\lfloor y\rfloor<k$ jumps $V$ and then is stuck, never leaving the small region $x<4$, $y<4$; it certainly stays within $2018$ of the origin.
\end{itemize}
Points with $x\ge4$ escape directly. So the set is the staircase region
\[
\{0\le x<1,\ 0\le y<4\}\cup\{1\le x<2,\ 0\le y<3\}\cup\{2\le x<3,\ 0\le y<2\}\cup\{3\le x<4,\ 0\le y<1\},
\]
shaded in the figure.

\begin{asy}
usepackage("amsmath");
size(6.5cm);
defaultpen(fontsize(10pt));
for(int i=0;i<=6;++i){draw((i,0)--(i,6.2),gray(0.88));draw((0,i)--(6.2,i),gray(0.88));}
path R=(0,0)--(4,0)--(4,1)--(3,1)--(3,2)--(2,2)--(2,3)--(1,3)--(1,4)--(0,4)--cycle;
fill(R,paleblue);draw(R,blue+linewidth(0.9));
draw((0,0)--(6.6,0),Arrow(TeXHead));draw((0,0)--(0,6.6),Arrow(TeXHead));
label("$x$",(6.6,0),E);label("$y$",(0,6.6),N);
for(int i=1;i<=4;++i){label("$"+string(i)+"$",(i,0),S);label("$"+string(i)+"$",(0,i),W);}
label("$0$",(0,0),SW);
pair P=(4.6,2.5);dot(P);
draw(P--(P+(-4,6)/3.2),red,Arrow(TeXHead));draw(P--(P+(1,-1)/1.6),red,Arrow(TeXHead));
label("$(-4,6)$",P+(-4,6)/3.2,W,red);label("$(1,-1)$",P+(1,-1)/1.6,E,red);
label("\mbox{}",(min(currentpicture,user=true)+max(currentpicture,user=true))/2,black); // keeps the text after the figure black when LaTeX typesets the labels inline
\end{asy}

Its area is $4+3+2+1=\boxed{10}$.
